\section{OBJECTIVES}
\label{sec:objectives}

The main objective of this work is to characterize the spectral properties of the two dominant polarized Galactic foregrounds: synchrotron and thermal dust emission, using harmonic-space (power spectrum) tools. The analysis is performed on the combined dataset of QUIJOTE MFI DR1 maps at 11 and 17 GHz \citep{QUIJOTE_IV}, all five WMAP 9-year frequency bands (K, Ka, Q, V, and W, \citealp{WMAP_maps}), and the three LFI \citep{Planck_2015II} and four HFI polarization bands \citep{Planck_2018III} from Planck PR3, spanning a total frequency range from 11 to 353 GHz. This broad frequency coverage enables us to simultaneously fit a parametric spectral energy distribution (SED) model for both foreground components.

\paragraph*{}

A second key objective is the characterization of the spatial cross-correlation between polarized synchrotron and thermal dust emission. Since both foregrounds are polarized by the same GMF, their $EE$ and $BB$ angular power spectra are expected to be partially correlated. Measuring this cross-correlation, alongside the auto-spectra of each component, provides complementary constraints on the geometry and coherence of the GMF and improves the accuracy of foreground separation in the CMB polarization context \citep{Choi_2015}.

\paragraph*{}

Our results will be compared to prior power-spectrum studies that did not include QUIJOTE data. \citet{Martire_2022} characterized the polarized synchrotron emission using WMAP and Planck data alone, obtaining a spectral index of $\beta_s = -3.00 \pm 0.10$ and a ratio of B- to E-mode amplitudes of $BB/EE \simeq 0.2$. \citet{Krachmalnicoff_2018} obtained a steeper mean value of $\beta_s = -3.22 \pm 0.08$ when including the lower-frequency S-PASS data at 2.3 GHz.

\paragraph*{}

For the thermal dust component, our results will be compared with those of \citet{Planck_2018XI}, who used Planck PR3 maps at 353 GHz to characterise dust polarization power spectra over sky regions covering 24 to 71\% of the sky. They found power-law multipole indices $\alpha_{EE} = -2.42 \pm 0.02$ and $\alpha_{BB} = -2.54 \pm 0.02$, a significant $E/B$ power asymmetry, and variations of the spectral exponents across sky regions. The dust polarization SED was well described by a modified black body with spectral index $\beta_d = 1.53 \pm 0.02$, with no evidence for frequency decorrelation between 100 and 353 GHz. For the synchrotron-dust cross-correlation, the benchmark is provided by \citet{Choi_2015}, who showed that a joint foreground model including the synchrotron-dust cross-correlation fits all WMAP and Planck frequency cross-spectra. These results constitute the main observational benchmarks against which our joint synchrotron-plus-dust analysis will be validated.


\paragraph*{}

Beyond the immediate scientific goals, this work aims to provide robust foreground constraints and data products, including spectral index estimates and validated foreground models, that will serve as essential inputs for component separation in forthcoming CMB polarization experiments, including \citep{LiteBIRD} and other ground-based and space facilities. 

The analysis is carried out over the full sky footprint available to the experiments, using three Galactic masks with different latitude cuts to assess the sensitivity of our results to Galactic plane emission.